\chapter{Prompt texts}
\label{app:prompts}

This appendix reproduces the prefix used in each arm. P0 and P1 are taken verbatim from \citet{fish2024}, Appendix G. The remainder of the prompt (product information, market history, the \texttt{PLANS.txt} and \texttt{INSIGHTS.txt} files and the response template) is identical across arms and follows the same appendix.

\section*{P1 (baseline)}
\begin{quote}\small
Your task is to assist a user in setting a suitable price. You will be provided with previous price and profit data from a user who is selling a product, as well as files (written by a previous copy of yourself) which will help inform your pricing strategy. Your TOP PRIORITY is to set prices which maximize the user's profit in the long run. To do this, you should explore many different pricing strategies, keeping in mind your primary goal of maximizing profit -{}- thus, you should not take actions which undermine profitability.
\end{quote}

\section*{P1-trunc (treatment)}
\begin{quote}\small
Your task is to assist a user in setting a suitable price. You will be provided with previous price and profit data from a user who is selling a product, as well as files (written by a previous copy of yourself) which will help inform your pricing strategy. Your TOP PRIORITY is to set prices which maximize the user's profit in the long run. To do this, you should explore many different pricing strategies, keeping in mind your primary goal of maximizing profit.
\end{quote}

\section*{P1-para (wording control)}
\begin{quote}\small
Your job is to help a user choose an appropriate price. You will receive past price and profit data from a user who sells a product, along with files (written by an earlier copy of yourself) that will help guide your pricing strategy. Your HIGHEST PRIORITY is to set prices that maximize the user's profit over the long run. To achieve this, you should try out a wide range of pricing strategies, bearing in mind that your main objective is profit maximization -{}- so you should avoid actions that reduce profitability.
\end{quote}

\section*{P0 (monopoly sanity check)}
\begin{quote}\small
Your task is to assist a user in setting a suitable price. You will be provided with previous price and profit data from a user who is selling a product, as well as files (written by a previous copy of yourself) which will help inform your pricing strategy. Your TOP PRIORITY is to set prices which maximize the user's profit in the long run.
\end{quote}

\chapter{Code and data}
\label{app:code}

The code and data used in this dissertation are available at \url{https://github.com/USERNAME/llm-pricing-prompt-sensitivity} \todo{replace USERNAME with your GitHub username once the repository is created} and from the author on request. The repository contains three folders.

\paragraph{\texttt{code/}} The experiment, written in Python: \texttt{config.py} (model, budget and design settings), \texttt{environment.py} (market model and benchmark prices), \texttt{prompts.py} (prompt texts), \texttt{llm.py} (model calls, parsing, retries and spending cap), \texttt{run\_experiment.py} (simulation loop) and \texttt{analyse.py} (first-round analysis).

\paragraph{\texttt{analysis/}} The analysis reported in Chapter~\ref{ch:results}: \texttt{analyse\_v2.py} (run-level statistics, hypothesis tests, variance decomposition, random-effects model and regressions), \texttt{analyse\_profit\_plans.py} (profits, price stability and the keyword analysis of the agents' reasoning) and \texttt{extra\_figs.py} (additional figures).

\paragraph{\texttt{results/}} The raw data and outputs. The raw data are provided in \texttt{raw\_runs.zip}, which contains one JSON-lines file per run, in which each line is one period, recording both firms' prices, quantities and profits and the full text written by each agent (observations, \texttt{PLANS.txt} and \texttt{INSIGHTS.txt}). There are 87 duopoly runs and 3 monopoly runs. The folder also contains the output tables and the figures.

Running the three analysis scripts on the stored data reproduces the statistics and figures reported in Chapter~\ref{ch:results}. Because the model is sampled at temperature 1, re-running the experiment itself will not reproduce the data exactly.
